\textbf{Skelný stav} je amorfní pevný stav bez pravidelné prostorové mřížky. Sklo je termodynamicky nestabilní a místo ostrého bodu tání postupně měkne.
\begin{itemize}
\item \textbf{Křemenné sklo:} síť tvoří SiO$_2$ bez modifikačních iontů. Má nízkou permitivitu ($\varepsilon_r\approx3{,}8$) a malé ztráty ($\tan\delta\approx2\cdot10^{-4}$); propouští viditelné i UV záření.
\item \textbf{Křemičité sklo:} síť obsahuje další oxidy a modifikační ionty. Permitivita i ztráty závisejí na složení; pohyblivé ionty zvyšují ztráty. Nízkoztrátová olovnatá skla mají $\varepsilon_r\approx7$--$8$, boritokřemičitá $4$--$5$. Křemičitá skla propouštějí především viditelné záření.
\end{itemize}
Pro \textbf{infračervené záření} se používají např. bezkyslíkatá chalkogenidová skla As--S.
